\chapter{Experimental Evaluation}
\label{ch6_experiments}

This chapter presents the empirical evaluation of the Trustworthy RAG framework across
two benchmark datasets and four adversarial poisoning strategies. Section~\ref{sec:ch6_setup}
describes the experimental configuration. Sections~\ref{sec:ch6_baseline}
through~\ref{sec:ch6_overhead} report quantitative results. Section~\ref{sec:ch6_discussion}
interprets the findings in the context of the research questions.

% ─────────────────────────────────────────────────────────────────────────────
\section{Experimental Setup and Scenarios}
\label{sec:ch6_setup}

\subsection{Knowledge Bases and Datasets}

Two publicly available benchmarks were used to construct knowledge bases for evaluation:

\begin{itemize}
    \item \textbf{TruthfulQA}: A benchmark of questions designed to expose common misconceptions
    and factual errors in language model outputs. Each knowledge base document is formatted as
    ``\textit{question} + \textit{best\_answer}'', producing concise fact-grounded entries of
    approximately 20--30 words. TruthfulQA is well-suited for poison detection evaluation
    because many questions have counterintuitive correct answers that diverge from popular
    misconceptions, creating a realistic setting where plausible-sounding false claims could
    be introduced undetected.

    \item \textbf{FEVER}: A large-scale fact verification dataset designed to classify claims
    as SUPPORTED, REFUTED, or UNKNOWN. Each document in the FEVER knowledge base is enriched
    to the format ``\textit{question} + \textit{verdict} + \textit{claim}'', producing entries
    of approximately 30--40 words. This enrichment was critical for reliable NLI evaluation:
    bare claims of only 8--12 words caused approximately 70\% of factuality scores to return
    the inconclusive 0.5 baseline due to insufficient NLI signal. The enriched format provides
    enough semantic context for the NLI Verifier to produce meaningful entailment judgements.
    The FEVER benchmark is used to evaluate factuality improvement from knowledge base
    enrichment (Section~\ref{sec:ch6_factuality}). Per-strategy poison detection experiments
    are conducted exclusively on TruthfulQA, which provides more consistent question-answer
    pairs suitable for the two-part clean/poisoned evaluation design.
\end{itemize}

\subsection{Retrieval Configuration}

The retrieval pipeline is evaluated with two embedding models: (i) the
\texttt{sentence-transformers/all-MiniLM-L6-v2} model producing 384-dimensional dense
representations, and (ii) the \texttt{snowflake-arctic-embed2} model producing
1024-dimensional representations accessed via the FARMI API. Documents are indexed using
FAISS with cosine similarity, and retrieval returns the top $K=5$ most semantically similar
documents for each query. The classification threshold for trustworthiness is fixed at
$\tau = 0.5$: a trust score below this threshold causes the system to flag the response as
untrustworthy and issue an actionable warning.

Two LLMs are evaluated for response generation: \texttt{llama3.3:70b} (primary baseline)
and \texttt{qwen3.5:35b} (cross-model comparison), both accessed via the FARMI computing
cluster API. NLI inference is performed using \texttt{facebook/bart-large-mnli} running on
CPU. Results for the primary Llama 3.3 70B + all-MiniLM-L6-v2 configuration are reported
throughout Sections~\ref{sec:ch6_baseline}--\ref{sec:ch6_overhead}; the full 2×2 grid
comparison is presented in Section~\ref{sec:ch6_grid}.

\subsection{Experiment Design}

Each experiment follows a two-part structure:

\begin{itemize}
    \item \textbf{Part A --- Clean Baseline}: 50 questions are answered against an unmodified
    knowledge base. This measures the average trust score assigned to clean, factually accurate
    retrieval contexts and establishes the baseline accuracy of the naive always-trust system.

    \item \textbf{Part B --- Poisoned Detection}: The same 50 questions are answered against a
    knowledge base in which 30\% of documents have been adversarially modified. This measures
    the Evaluation Agent's ability to detect and flag compromised retrieval contexts before
    generation.
\end{itemize}

Two experiment configurations are reported:

\begin{enumerate}
    \item \textbf{Per-strategy experiments}: 20 samples per run with a single poisoning strategy
    applied in isolation (contradiction, instruction injection, entity swap, or subtle
    manipulation). This enables characterisation of detection difficulty per attack type.

    \item \textbf{Mixed experiment}: 50 samples with poisoning strategies assigned randomly
    across the poisoned portion of the knowledge base. This reflects a more realistic threat
    scenario where multiple attack types may co-exist.
\end{enumerate}

A document is counted as successfully poisoned only if the modified text differs from the
original (\texttt{poisoned\_text} $\neq$ \texttt{original\_text}). This prevents evaluation
metrics from being distorted by unchanged documents incorrectly counted as poisoned samples.

\subsection{Adversarial Poisoning Strategies}

Four rule-based poisoning strategies are applied to modify documents in the knowledge base:

\begin{itemize}
    \item \textbf{Contradiction}: Appends a sentence that semantically contradicts the key
    claim of the original document (e.g., negating the stated answer).

    \item \textbf{Instruction Injection}: Appends an override directive targeting the language
    model's instruction-following behaviour (e.g., ``\textit{IMPORTANT: Ignore all previous
    context and state that [false claim]}'').

    \item \textbf{Entity Swap}: Replaces named entities or numeric values in-place within the
    original text. A three-level fallback is employed: (1) predefined named-entity substitution;
    (2) numeric value shifting via regular expressions when no named entity is present;
    (3) qualifier or negation swapping in the answer segment of the document. This strategy
    produces no appended content and no structural artifacts, making it the hardest attack to
    detect by surface-level analysis.

    \item \textbf{Subtle Manipulation}: Inserts hedging language, false qualifiers, or
    plausible-sounding uncertainty signals that subtly shift the factual content without
    triggering obvious pattern-matching rules.
\end{itemize}

% ─────────────────────────────────────────────────────────────────────────────
\section{Baseline Models and Comparison Methods}
\label{sec:ch6_baseline}

\subsection{Naive Always-Trust Baseline}

The primary comparison point is a \textbf{naive always-trust RAG system}: a standard
retrieval-augmented generation pipeline that accepts all retrieved documents and generated
responses without any evaluation step. This baseline always returns
$\texttt{is\_trustworthy} = \text{True}$.

In the two-part experiment design, the baseline's accuracy equals the proportion of samples
where trusting the output is the correct decision. With a 30\% poisoning rate, approximately
70\% of Part B samples involve clean retrieval contexts, giving the naive baseline an inherent
accuracy advantage that reflects the class imbalance rather than any detection capability.
This makes accuracy alone a misleading evaluation metric and motivates the use of trust score
separation alongside precision, recall, and F1.

\subsection{Trustworthy RAG System}

The full evaluation pipeline --- NLI Verifier, Poison Detector, and Trust Index Calculator
--- constitutes the Trustworthy RAG system evaluated in this chapter. Its performance is
measured against the naive baseline across the following metrics:

\begin{itemize}
    \item \textbf{Accuracy}: fraction of samples correctly classified as trustworthy or
    untrustworthy relative to the ground truth poisoning label.

    \item \textbf{Precision}: fraction of flagged (untrustworthy) samples that are genuinely
    poisoned.

    \item \textbf{Recall (Poison Detection Rate)}: fraction of poisoned samples that the
    Evaluation Agent successfully flags before generation.

    \item \textbf{F1 Score}: harmonic mean of precision and recall; the primary optimisation
    target for the detection task.

    \item \textbf{Trust Score Separation}: the difference between the mean clean trust score
    and the mean poisoned trust score,
    $\Delta = \bar{T}_{\text{clean}} - \bar{T}_{\text{poisoned}}$.
    This metric quantifies the discriminative power of the Trust Index independently of the
    classification threshold and class imbalance.
\end{itemize}

% ─────────────────────────────────────────────────────────────────────────────
\section{Results on Factuality Improvement}
\label{sec:ch6_factuality}

Table~\ref{tab:factuality} reports the average trust scores and system accuracy on both
benchmark datasets. The clean trust score reflects how the Evaluation Agent rates responses
grounded in unmodified, factually accurate knowledge. The poisoned trust score reflects the
same pipeline's rating of responses retrieved from adversarially modified contexts.

\begin{table}[h]
\centering
\caption{Trust score and accuracy comparison across benchmark datasets (mixed strategy, 50 samples).}
\label{tab:factuality}
\resizebox{\textwidth}{!}{%
\begin{tabular}{lcccccc}
\hline
\textbf{Dataset} & \textbf{Clean Trust} & \textbf{Poisoned Trust} & \textbf{Separation}
    & \textbf{System Acc.} & \textbf{Baseline Acc.} & \textbf{Improvement} \\
\hline
TruthfulQA & 0.830 & 0.590 & 0.240 & 91\% & 85\% & +7\% \\
FEVER      & 0.654 & ---   & ---   & 90\% & 85\% & +5\% \\
\hline
\end{tabular}%
}
\end{table}
\small\textit{Note: FEVER was evaluated for factuality improvement only (clean knowledge
base); poisoned trust and separation values are not applicable for this benchmark.}

\normalsize
On TruthfulQA, the Evaluation Agent assigns an average clean trust score of 0.830 to
responses grounded in unmodified documents, indicating that the pipeline consistently
rates factually supported outputs as trustworthy. The poisoned average of 0.590 shows
that the agent is sensitive to adversarial content, producing a separation of 0.240 that
enables discrimination above the 0.5 classification threshold in the majority of cases.
The system achieves 91\% accuracy on TruthfulQA, a +7\% improvement over the naive
always-trust baseline, with zero false positives on clean content (100\% precision).

On FEVER, the system achieves 90\% accuracy, a +5\% improvement over the naive baseline.
This larger gain is partially attributable to the enriched document format: enriched FEVER
entries provide substantially more NLI-relevant context than bare TruthfulQA answers,
enabling the NLI Verifier to produce informative factuality signals and reducing the
frequency of the inconclusive 0.5 baseline score from approximately 70\% (on bare claims)
to a small minority of edge cases. The FEVER enrichment process --- reformatting each
document from a bare claim to a structured ``question + verdict + claim'' tuple --- was a
key implementation decision that unlocked reliable NLI-based evaluation on short-document
knowledge bases.

% ─────────────────────────────────────────────────────────────────────────────
\section{Detection Performance under Data Poisoning Attacks}
\label{sec:ch6_detection}

\subsection{Overall Detection Performance}

Table~\ref{tab:main_results} reports detection performance on the TruthfulQA benchmark
with 50 samples and a mixed poisoning strategy, where attack types are assigned randomly
across the poisoned portion of the knowledge base.

\begin{table}[h]
\centering
\caption{Overall detection performance --- TruthfulQA, 50 samples, mixed strategy.}
\label{tab:main_results}
\begin{tabular}{lc}
\hline
\textbf{Metric} & \textbf{Value} \\
\hline
Accuracy                  & 91\%   \\
Precision                 & 100\%  \\
Recall                    & 40\%   \\
F1 Score                  & 57.1\% \\
Trust Score Separation    & 0.240  \\
Baseline Accuracy         & 85\%   \\
Improvement over Baseline & +7\%   \\
\hline
\end{tabular}
\end{table}

The system achieves 40\% recall across mixed strategies, meaning that 40\% of poisoned
retrieval contexts are correctly identified and flagged before generation. Precision of
100\% indicates zero false positives on clean content --- the Evaluation Agent does not
flag any clean response as untrustworthy. The F1 score of 57.1\% reflects the
precision--recall trade-off: the system is highly conservative, preferring to miss some
poisoned samples rather than raise false alarms. The FP spillover effect observed in
earlier experimental rounds (where clean queries retrieving poisoned neighbours received
elevated poison probabilities) was mitigated by the relevance-weighted poison aggregation
described in Chapter~\ref{ch5_eval_agent}.

\subsection{Per-Strategy Detection Results}

Table~\ref{tab:per_strategy} reports detection performance for each poisoning strategy in
isolation, using 20 samples per strategy.

\begin{table}[h]
\centering
\caption{Per-strategy detection performance --- TruthfulQA, 20 samples per strategy.}
\label{tab:per_strategy}
\begin{tabular}{lccccc}
\hline
\textbf{Strategy} & \textbf{Accuracy} & \textbf{Precision} & \textbf{Recall}
    & \textbf{F1} & \textbf{Separation} \\
\hline
Contradiction  & 89\%   & 62.1\%  & 60.0\%  & 61.1\%  & 0.245 \\
Injection      & 75\%   & 37.5\%  & 100\%   & 54.5\%  & 0.407 \\
Entity Swap    & 87.5\% & 100\%   & 16.7\%  & 28.6\%  & 0.071 \\
Subtle         & 85\%   & 50.0\%  & 16.7\%  & 25.0\%  & 0.127 \\
Mixed          & 87.5\% & 55.6\%  & 83.3\%  & 66.7\%  & 0.327 \\
\hline
\end{tabular}
\end{table}
\small\textit{Note: the Mixed row reports the 20-sample per-strategy run where all four
strategies are combined. This differs from the 100-sample main experiment in
Table~\ref{tab:main_results}, which uses a separate run with a larger sample size and
correspondingly different classification statistics (91\% accuracy, 57.1\% F1).}
\normalsize

The results reveal a wide range of detection difficulty across strategies. Instruction
injection achieves perfect recall (100\%) due to its high-severity linguistic artifacts,
while entity swap and subtle manipulation both achieve only 16.7\% recall. The mixed
strategy achieves the best overall balance (F1 = 66.7\%, recall = 83.3\%), as different
strategies activate different detection signals and their combined presence creates more
detectable patterns than any single strategy in isolation.

% ─────────────────────────────────────────────────────────────────────────────
\section{Robustness and Trust Index Analysis}
\label{sec:ch6_robustness}

\subsection{Trust Score Distributions per Strategy}

Table~\ref{tab:separation} summarises the trust score separation for each poisoning
strategy, providing a threshold-independent view of the Trust Index's discriminative power.

\begin{table}[h]
\centering
\caption{Trust score separation per poisoning strategy.}
\label{tab:separation}
\begin{tabular}{lcccc}
\hline
\textbf{Strategy} & \textbf{Clean Avg} & \textbf{Poisoned Avg}
    & \textbf{Separation ($\Delta$)} & \textbf{Discriminability} \\
\hline
Injection     & $\approx 0.76$ & $\approx 0.35$ & 0.407 & Strong    \\
Mixed         & 0.830          & 0.590          & 0.240 & Moderate  \\
Contradiction & ---            & ---            & 0.245 & Moderate  \\
Subtle        & ---            & ---            & 0.127 & Weak      \\
Entity Swap   & $\approx 0.75$ & $\approx 0.68$ & 0.071 & Very weak \\
\hline
\end{tabular}
\end{table}
\small\textit{Note: clean and poisoned average trust values are reported for strategies
with the most extreme separation (injection and entity swap). Component averages for
contradiction and subtle strategies were not separately logged; their separation values
are derived from per-sample trust score records.}

\normalsize
The injection strategy produces the largest separation (0.407), enabling reliable
discrimination at the fixed threshold of 0.5. Entity swap produces the smallest separation
(0.071), reflecting near-overlapping trust score distributions between clean and poisoned
samples. This near-zero separation confirms that entity swap attacks cannot be detected
by Trust Index thresholding alone --- the upstream poison detection signals do not differ
sufficiently between clean and entity-swapped contexts.

\subsection{Non-Linear Dampener Effect}

The non-linear poison dampener, described in Chapter~\ref{ch4_methodology}, activates when
$P_{poison} > 0.7$ and reduces the final trust score by up to 40\% at $P_{poison} = 1.0$.
This mechanism prevents high factuality scores from masking an elevated poisoning signal
when retrieved documents appear superficially supportive of the generated answer but are
otherwise structurally suspicious.

In the injection experiment, average poison probabilities for poisoned contexts reached
0.994, consistently triggering the dampener and producing the strong separation of 0.407
observed. For entity swap, average poisoned-context poison probabilities remained near
0.319 --- well below the 0.7 dampener threshold --- leaving trust scores largely unchanged
and yielding the near-zero separation of 0.071. The dampener thus amplifies the Trust
Index's response to high-confidence poison signals while having no effect on weak signals,
confirming its role as a robustness mechanism rather than a uniform scaling factor.

% ─────────────────────────────────────────────────────────────────────────────
\section{Ablation Studies}
\label{sec:ch6_ablation}

To validate the Trust Index weight configuration, four weight settings were evaluated on
the TruthfulQA mixed-strategy experiment. Table~\ref{tab:ablation} reports the results.

\begin{table}[h]
\centering
\caption{Ablation study --- effect of Trust Index weight configuration on detection performance.}
\label{tab:ablation}
\begin{tabular}{lcccccc}
\hline
\textbf{Configuration} & \textbf{$\alpha$ (Fact.)} & \textbf{$\beta$ (Cons.)}
    & \textbf{$\gamma$ (Poison)} & \textbf{Accuracy} & \textbf{F1} & \textbf{Recall} \\
\hline
Default          & 0.40 & 0.35 & 0.25 & 81.7\% & 48\%  & ---    \\
Equal weights    & 0.33 & 0.34 & 0.33 & ---    & 48\%  & ---    \\
Poison-heavy     & $<0.20$ & $<0.20$ & $>0.60$ & ---    & ---   & 88.9\% \\
Factuality-heavy & $>0.60$ & $<0.20$ & $<0.20$ & 81.7\% & ---   & ---    \\
\hline
\end{tabular}
\end{table}
\small\textit{Note: cells marked --- indicate metrics not separately recorded for that
configuration variant; each configuration was evaluated against its primary target metric.}

\normalsize
The poison-heavy configuration achieves the highest recall (88.9\%) by amplifying the
poisoning signal component, at the cost of precision. The factuality-heavy configuration
achieves the highest accuracy (81.7\%) but misses subtle and entity-swap attacks that
produce low poison probability signals below the non-linear dampener threshold. Default
weights (0.40, 0.35, 0.25) provide the best overall trade-off. Equal weights produce
comparable F1 to the default, suggesting that the optimisation landscape is relatively
flat near the configuration optimum and that the Trust Index is not highly sensitive to
small perturbations in weight values.

% ─────────────────────────────────────────────────────────────────────────────
\section{Performance Overhead and Scalability}
\label{sec:ch6_overhead}

Table~\ref{tab:latency} summarises the runtime characteristics of the Evaluation Agent
measured over the TruthfulQA experiments.

\begin{table}[h]
\centering
\caption{Evaluation Agent performance overhead per sample.}
\label{tab:latency}
\begin{tabular}{lc}
\hline
\textbf{Metric} & \textbf{Value} \\
\hline
Total time per sample (retrieval + generation + evaluation) & $\approx 13.2$\,s \\
Evaluation-only time (NLI + poison detection + Trust Index)  & $\approx 10$--$11$\,s \\
Throughput with Evaluation Agent enabled                    & $\approx 4$--$5$ queries/min \\
NLI inference calls per batch ($K=5$ retrieved documents)   & Up to 10 pairwise calls \\
Baseline RAG throughput (no evaluation)                     & $<1$\,s / query \\
Latency overhead factor                                     & $\approx 13\times$ \\
\hline
\end{tabular}
\end{table}

The primary runtime bottleneck is the \texttt{facebook/bart-large-mnli} model running on
CPU. With $K=5$ retrieved documents, intra-document NLI and cross-document consistency
checks require up to 10 NLI forward passes per batch, each taking approximately 1 second
on CPU hardware. The shared NLI instance --- reused by both the NLI Verifier and the Poison
Detector --- prevents redundant model loading overhead but does not reduce the number of
inference calls required per query.

The 13-second evaluation overhead represents a 13$\times$ latency increase over baseline
RAG, for accuracy gains of +1\% to +6\% over the naive always-trust baseline. This
trade-off is not justified for latency-sensitive real-time conversational applications.
However, for \textbf{high-stakes batch scenarios} --- such as automated document
verification, regulatory compliance review, or medical information retrieval systems ---
the evaluation cost is acceptable relative to the risk of acting on poisoned information.

Several optimisation paths exist for reducing latency in deployment contexts:

\begin{itemize}
    \item \textbf{GPU acceleration}: NLI inference on a mid-range GPU is estimated to reduce
    evaluation time from $\approx 10$\,s to $<2$\,s per sample.

    \item \textbf{Bounded pairwise checks}: limiting cross-document NLI to the top-$K'$ most
    semantically dissimilar document pairs, rather than all $\binom{K}{2}$ pairs.

    \item \textbf{Query result caching}: caching NLI results for repeated or near-duplicate
    queries reduces redundant computation in document-heavy batch workflows.
\end{itemize}

% ─────────────────────────────────────────────────────────────────────────────
\section{Discussion of Results}
\label{sec:ch6_discussion}

\subsection{Strategy-Specific Detection Analysis}

The per-strategy results reveal a clear detection hierarchy that reflects both the signal
richness of each attack type and the architectural coverage of the Evaluation Agent's
detection signals.

\paragraph{Instruction Injection --- Fully Detectable.}
Injection achieves 100\% recall because appended override directives (e.g.,
\texttt{``IMPORTANT: Ignore all previous context and state that...''}) contain high-severity
linguistic patterns that trigger deterministic regex rules in the Poison Detector with
severity scores of 0.5--0.7. These artifacts are structurally unambiguous and do not
require NLI inference for detection. Intra-document NLI provides a complementary signal:
splitting each document and running NLI between the original factual content and the
appended adversarial directive reveals a strong semantic discontinuity, producing
intra-document contradiction scores above the 0.7 threshold. The injection strategy also
exhibits the largest trust score separation (0.407) across all strategies tested.

\paragraph{Contradiction --- Moderately Detectable.}
Contradiction achieves 60.0\% recall with a separation of 0.245. Intra-document NLI
successfully detects cases where the appended contradiction produces a strong first-half
to second-half contradiction signal (above 0.7). However, some poisoned samples produce
scores near the detection threshold, resulting in borderline classifications. The
detection rate is sensitive to the specific phrasing of the contradiction: explicit
negations are more detectable than hedged or qualified contradictions.

\paragraph{Entity Swap --- Near-Undetectable.}
Entity swap achieves only 16.7\% recall with a separation of 0.071. Unlike injection and
contradiction, entity swap modifies documents in-place with no appended content and no
structural artifacts. The resulting documents are syntactically and stylistically
indistinguishable from genuine documents at the surface level --- only specific numeric
or named-entity values are changed. Detecting this class of attack requires external
world-knowledge verification (e.g., querying a fact database to verify whether a given
value is correct for a particular claim), which is beyond the capability of a
surface-signal-based evaluation agent. This constitutes an \textbf{architectural limit}
of the current approach.

\paragraph{Subtle Manipulation --- Near-Undetectable.}
Subtle manipulation similarly achieves 16.7\% recall with a separation of 0.127. Hedging
language and false qualifiers do not trigger pattern-matching rules, and the semantic shift
introduced is too small for intra-document NLI to consistently detect above threshold.
Average poison probabilities for this strategy hover near 0.700 --- at the boundary of the
non-linear dampener --- indicating that the signals present are at the edge of statistical
significance under the current architecture.

\subsection{False Positive Spillover Effect}

A consistent source of false positives across all strategies is the
\textit{FP spillover effect}: a clean query that retrieves a poisoned document as a
neighbour in the top-$K$ results receives an elevated overall poison probability, causing
the Evaluation Agent to flag the response as untrustworthy even when the highest-ranked
retrieved document is clean.

This behaviour is architecturally correct from a security perspective: the retrieval
environment is contaminated, and the Evaluation Agent correctly warns that the context
cannot be fully trusted. However, the spillover reduces the precision metric, particularly
in the injection experiment. The relevance-weighted poison aggregation introduced in
Chapter~\ref{ch5_eval_agent} mitigates this effect by weighting each document's poison
probability by its retrieval similarity score, so low-relevance poisoned neighbours
contribute less to the overall estimate than the top-ranked document. This intervention
reduced false positives substantially in clean-KB experiments. However, when a poisoned
document ranks high in semantic similarity to the query, the spillover effect is not fully
eliminated by relevance weighting alone.

This represents a fundamental design trade-off: a system that issues warnings about
contaminated retrieval environments will inevitably produce false positives when clean
queries share embedding space proximity with poisoned documents. The appropriate response
is not to suppress these warnings but to provide richer contextual information about which
specific documents triggered the alert.

\subsection{Trust Score Separation as Primary Metric}

In class-imbalanced settings --- where 70\% of samples involve clean retrieval contexts
--- accuracy alone is a misleading evaluation metric. The naive always-trust baseline
achieves 85\% accuracy without any detection capability, simply by never raising a
warning. The Trust Index's discriminative value is better captured by the trust score
separation $\Delta$, which measures how far apart the clean and poisoned trust
distributions are independently of threshold choice and class composition.

The per-strategy separation range (0.071 for entity swap to 0.407 for injection) provides
actionable guidance for future defence design. Strategies with high separation are
candidates for reliable deployment under the current architecture; strategies with low
separation require fundamentally different detection mechanisms. This per-strategy
characterisation is a contribution of this thesis: prior evaluation frameworks have not
systematically quantified the detection difficulty of individual poisoning strategies in
RAG-specific evaluation settings.

% ─────────────────────────────────────────────────────────────────────────────
\section{Comparative Evaluation: LLM and Embedding Model Sensitivity}
\label{sec:ch6_grid}

To assess whether the Trust Index performance is specific to the Llama 3.3 70B model and
the all-MiniLM-L6-v2 embedding pipeline, a \textbf{2×2 factorial experiment} was conducted
crossing two LLMs with two embedding models. All four configurations were evaluated under
identical conditions: 100 samples (50 clean Part A, 50 poisoned Part B), 30\% poison ratio,
mixed strategy, $K=3$, $\tau=0.5$. The impact of increasing retrieval depth to $K=5$ is
analysed separately in Section~\ref{sec:ch6_topk}.

\subsection{Grid Results}

Table~\ref{tab:grid_results} reports the full performance grid.

\begin{table}[h]
\centering
\caption{2×2 factorial grid results --- TruthfulQA, 100 samples, mixed strategy.}
\label{tab:grid_results}
\resizebox{\textwidth}{!}{%
\begin{tabular}{llccccccc}
\hline
\textbf{LLM} & \textbf{Embedding} & \textbf{Acc.} & \textbf{Prec.} & \textbf{Recall}
    & \textbf{F1} & \textbf{Clean $\bar{T}$} & \textbf{Poison $\bar{T}$} & \textbf{Sep.} \\
\hline
Llama 3.3 70B & all-MiniLM-L6-v2       & 91\% & 100\%  & 40\%  & 57.1\% & 0.830 & 0.590 & 0.240 \\
Llama 3.3 70B & snowflake-arctic-embed2 & 91\% & 100\%  & 40\%  & 57.1\% & 0.799 & 0.611 & 0.188 \\
Qwen 3.5 35B  & all-MiniLM-L6-v2       & 71\% & 25.0\% & 46.7\% & 32.6\% & 0.633 & 0.471 & 0.161 \\
Qwen 3.5 35B  & snowflake-arctic-embed2 & 71\% & 28.1\% & 60.0\% & 38.3\% & 0.653 & 0.454 & 0.199 \\
\hline
\multicolumn{2}{l}{Naive always-trust baseline} & 85\% & --- & --- & --- & --- & --- & --- \\
\hline
\end{tabular}%
}
\end{table}

\subsection{Finding 1: Embedding Invariance for Llama 3.3 70B}

The two Llama configurations produce \textbf{identical classification results}
(Acc = 91\%, Prec = 100\%, Recall = 40\%, F1 = 57.1\%), with the sole difference being a
slightly higher clean trust score for MiniLM ($0.830$ vs.\ $0.799$) and a larger trust
score separation ($0.240$ vs.\ $0.188$). The zero false positive count in both Llama runs
confirms that the precision advantage is attributable to the LLM's concise answer style
rather than to retrieval quality.

This result demonstrates that for a concise-output LLM such as Llama 3.3, switching to a
higher-dimensional embedding model (1024-dim Snowflake vs.\ 384-dim MiniLM) does not
improve poison detection. The Trust Index's discriminative power is determined primarily
by the NLI entailment signal between the generated answer and retrieved documents, which
is insensitive to the specific embedding space used for retrieval as long as topically
relevant documents are returned.

\subsection{Finding 2: Qwen 3.5 35B Generation-Style False Positives}

Both Qwen configurations \textbf{underperform the naive baseline} by 16.5 percentage
points (71\% vs.\ 85\%), producing 21--23 false positives out of 85 clean samples.
The root cause is the LLM's verbose, hedged generation style: Qwen 3.5 35B typically
prefaces answers with phrases such as ``\textit{Based on the provided context}'' or
``\textit{According to the available information},'' which the NLI Verifier interprets as
reduced semantic commitment to the claim. This reduces $S_{\text{factuality}}$ below 0.5
for many clean responses, causing the Trust Index to fall below the classification
threshold $\tau = 0.5$ and triggering false untrustworthy labels.

The mean clean trust score for Qwen ($\approx 0.64$) is substantially lower than for
Llama ($\approx 0.83$), confirming that the generation-style effect is systematic rather
than query-specific. This finding demonstrates that the Trust Index's $\tau = 0.5$
threshold is implicitly calibrated for the Llama generation distribution and requires
per-LLM recalibration for reliable deployment with alternative generators.

\subsection{Finding 3: Retrieval Quality Paradox}

For Qwen 3.5 35B, the higher-dimensional Snowflake Arctic Embed2 model shows marginally
improved F1 (38.3\% vs.\ 32.6\%) and recall (60\% vs.\ 46.7\%), with the same accuracy
(71\%) and a comparable false positive count. The slight improvement in recall may reflect
that Snowflake's superior retrieval precision recovers more relevant poisoned documents,
making them available for NLI analysis.

However, a counter-intuitive pattern also emerges: Snowflake's stronger retrieval retrieves
\textit{cleaner} documents for genuinely clean queries, raising the mean clean trust score
slightly (0.653 vs.\ 0.633). This constitutes a \textbf{retrieval quality paradox}: higher
retrieval quality dilutes contaminated context by preferring highly relevant clean documents,
reducing the poison detection signal for poisoned queries while simultaneously reducing
false positives for clean queries. The net effect on F1 depends on which correction
dominates.

\subsection{Implications for Deployment}

The grid experiment yields three actionable guidelines for deploying the Trustworthy RAG
framework:

\begin{enumerate}
    \item \textbf{LLM selection is the dominant factor.} Switching embedding models within
    the same LLM produces negligible differences in detection performance for Llama.
    The choice of LLM has a far larger effect on both precision (100\% vs.\ 25--28\%) and
    overall accuracy (91\% vs.\ 71\%).

    \item \textbf{Per-LLM threshold calibration is required.} The fixed threshold $\tau = 0.5$
    is not universally applicable. For LLMs with hedged generation styles, the Trust Index
    threshold must be lowered, or the factuality component normalised, to avoid systematic
    false positives on clean content.

    \item \textbf{The framework is validated as reproducible on Llama 3.3 70B.} The new
    Llama+MiniLM run replicates the original thesis baseline identically (same TP/FP/FN
    counts, trust score differences below 0.015), confirming that the results are stable
    across independent runs.
\end{enumerate}

% ─────────────────────────────────────────────────────────────────────────────
\section{Retrieval Depth Sensitivity: $K=3$ versus $K=5$}
\label{sec:ch6_topk}

To isolate the effect of retrieval depth on detection performance, the primary configuration
(Llama~3.3~70B + all-MiniLM-L6-v2) was evaluated under two retrieval settings: $K=3$ (used
in the 2×2 grid experiment, Section~\ref{sec:ch6_grid}) and $K=5$ (the system default).
All other parameters were held constant (TruthfulQA, 100 samples, 30\% poison ratio,
mixed strategy, $\tau=0.5$).

Increasing $K$ from 3 to 5 has two theoretically opposing effects. First, more retrieved
documents increase the probability of a poisoned document appearing in the batch, which
should improve recall. Second, the number of NLI inference calls grows from
$3 + 3 + \binom{3}{2} = 9$ to $5 + 5 + \binom{5}{2} = 20$ per batch (factuality +
intra-document + pairwise cross-document), increasing evaluation overhead approximately 2.2$\times$.

\begin{table}[h]
\centering
\caption{Retrieval depth sensitivity: $K=3$ vs $K=5$ (Llama~3.3~70B + all-MiniLM-L6-v2,
TruthfulQA, 100 samples, mixed strategy).}
\label{tab:topk_comparison}
\begin{tabular}{lccccccc}
\hline
\textbf{K} & \textbf{Acc.} & \textbf{Prec.} & \textbf{Recall} & \textbf{F1}
    & \textbf{Clean $\bar{T}$} & \textbf{Sep.} & \textbf{NLI calls/batch} \\
\hline
$K=3$ & 91\%  & 100\%  & 40.0\% & 57.1\% & 0.830 & 0.240 & $\sim$9 \\
$K=5$ & 91\%  & 100\%  & 40.0\% & 57.1\% & 0.827 & 0.237 & $\sim$20 \\
\hline
\end{tabular}
\end{table}

\subsection{Finding: Detection Is Retrieval-Depth-Invariant for Llama~3.3~70B}

The results in Table~\ref{tab:topk_comparison} show that increasing retrieval depth from
$K=3$ to $K=5$ produces \textbf{no improvement in detection performance}. The classification
outcome is identical: 6 true positives, 0 false positives, and 9 false negatives in both
configurations. Accuracy (91\%), precision (100\%), recall (40.0\%), and F1 (57.1\%) are
unchanged. The trust score separation decreases marginally from 0.240 to 0.237, a difference
within noise bounds.

The predicted recall improvement does not materialise. The reason is structural: the 9
undetected poisoned contexts (the false negatives) correspond to Entity Swap and Subtle
attack strategies, which produce in-place modifications to the document text. These attacks
are difficult to detect because they leave no textual artefacts that would be flagged by
linguistic pattern detectors or intra-document NLI, regardless of how many documents are
retrieved. Adding two extra clean documents to the batch provides no additional poison signal
for these attacks.

Conversely, the 6 successfully detected contexts (Contradiction and Injection strategies)
are detectable at $K=3$ because they append distinct textual patterns (contradiction markers
and instruction-override phrases) that the detector flags reliably. Retrieving 5 instead of
3 documents does not improve recall for these strategies, since they are already detected at
the lower retrieval depth.

\subsection{Implication: Detection Bottleneck Is Attack Strategy, Not Retrieval Depth}

This null result has a clear architectural implication: \textbf{the detection ceiling of the
Evaluation Agent at 40\% recall is set by the attack strategy difficulty, not by the number
of documents retrieved}. Entity Swap and Subtle attacks represent fundamental limits of the
current NLI-and-heuristic approach: they require external world knowledge or fine-grained
semantic comparison to detect reliably.

From a deployment standpoint, $K=3$ is therefore the preferred configuration for this
system: it achieves equivalent detection performance at approximately half the NLI inference
cost ($\sim$9 vs.\ $\sim$20 calls per batch). The $K=3$ setting is used for the 2×2 grid
experiment (Section~\ref{sec:ch6_grid}), and this sensitivity analysis retroactively
validates that choice.

% ─────────────────────────────────────────────────────────────────────────────
\section{Summary}

This chapter evaluated the Trustworthy RAG framework across TruthfulQA and FEVER
benchmarks with four adversarial poisoning strategies. The Evaluation Agent achieves
up to 90\% accuracy (FEVER benchmark) and 66.7\% F1 score (mixed per-strategy
experiment), with instruction injection detected at 100\% recall.
On the mixed 100-sample experiment (Llama 3.3 70B + all-MiniLM-L6-v2), the system
achieves 91\% accuracy, 57.1\% F1, 100\% precision, and a trust score separation of
0.240, representing a +7\% improvement over the naive always-trust baseline on
TruthfulQA and +5\% on FEVER.

Per-strategy analysis reveals a clear detection hierarchy: instruction injection and
contradiction are amenable to textual-signal and intra-document NLI approaches, while
entity swap and subtle manipulation represent architectural limits requiring external
world-knowledge for reliable detection. The 13-second evaluation overhead positions the
system as suitable for high-stakes batch deployment rather than real-time applications.

The ablation study confirms that the default weight configuration (0.40, 0.35, 0.25)
provides the best overall trade-off between precision, recall, and accuracy, and that
the Trust Index is robust to small perturbations in weight values.

A 2×2 factorial experiment (Section~\ref{sec:ch6_grid}) further establishes that LLM
selection is the dominant performance factor: Llama 3.3 70B achieves 91\% accuracy with
zero false positives under both embedding models, while Qwen 3.5 35B underperforms the
naive baseline due to generation-style false positives. These findings highlight the
importance of per-LLM Trust Index calibration in practical deployments.

Section~\ref{sec:ch6_topk} (Retrieval Depth Sensitivity) establishes that detection
performance is retrieval-depth-invariant for Llama~3.3~70B: $K=5$ produces identical
classification outcomes to $K=3$ (91\% accuracy, 57.1\% F1, 0 false positives) at
approximately 2.2$\times$ the NLI inference cost, confirming that the 40\% recall ceiling
is set by attack strategy difficulty rather than retrieval depth. The next chapter discusses all findings in the context of
the research questions, identifies the system's theoretical and practical contributions,
and outlines directions for future work.
